\section{Introduction}
\label{sec:introduction}

Bangladeshi agencies publish environmental measurements that can span several
decades. Weather stations record temperature, rain gauges record rainfall, and
river monitoring points sample water for oxygen and acidity. The implemented
records cover the period from \mYearMin{} to \mYearMax.

These publications are difficult to use directly in a database. The sheets
support human reading, so they contain merged headings, repeated date columns,
station names placed outside the data rows, and source-specific ways of marking
missing readings. A person can infer these meanings from the page layout, but
an extraction program needs them to be stated consistently.

A person reads this without noticing the difficulty, a query cannot. The
question of what the mean maximum temperature in Bogura was in April 2019
requires a machine to recognize alternate station spellings, identify the
correct date column, and distinguish a missing reading from zero. Reliable
queries depend on explicit rules for each of these cases.

This project makes the measurements queryable. The group designs a relational
schema, extracts the published material into it, normalizes the result to
Boyce--Codd Normal Form and records what the process costs.

\subsection{Background and Motivation}
\label{subsec:bm}

The work belongs to MODELBD, which the cohort develops across seven sectors
under a Scrum of Scrums arrangement \cite{scrumguide2020}. Group-07 holds the
environment sector at Scrum Level-01.

The motivation is a gap between publication and use. Bangladesh is highly
exposed to climate-related hazards and its own agencies hold long-run evidence
needed to study that exposure \cite{worldbank_climate_bd}. A researcher who
wants to compare rainfall in the 1970s against rainfall in the 2010s currently
opens several spreadsheets and reconciles them by hand. That reconciliation is
slow and is not reproducible when researchers resolve names such as Bogura in
different ways.

A normalized database removes that variance. Once the organisations' figures
occupy one schema with one spelling per station and one notation for a missing
value, the comparison becomes a query and the query returns the same answer to
everyone who runs it.

\subsection{Problem Statement}
\label{subsec:ps}

Three recurring problems shape the work.

\begin{description}[leftmargin=1.15cm,style=nextline]
  \item[P1] The published form is not relational. Column groups repeat, headers
  merge across cells and a row's identity can depend on a label above it. A
  sheet in this form has no dependable primary key and violates first normal
  form before any design decision is made.
  \item[P2] The same real entity carries several names. Station and place names
  differ in spelling, capitalization and abbreviation between publications.
  Missing values are also written as symbols, spreadsheet errors, words and
  empty cells.
  \item[P3] Not every candidate publication belongs in the final model. Some
  material is narrative, outside the approved scope or lacks a reliable join
  key. Even retained files use different grains and require a reproducible path
  from extraction to a tested database.
\end{description}

\subsection{Project Objectives}
\label{subsec:objectives}

The project has three report-level objectives that respond to these problems.

\begin{description}[leftmargin=1.25cm,style=nextline]
  \item[OB1] Collect, verify, extract and prepare environmental records from
  published sources while retaining source traceability and data-quality
  decisions.
  \item[OB2] Design an integrated environmental data model, normalize it from
  first normal form to BCNF and implement the approved relations in SQLite.
  \item[OB3] Demonstrate that the database is reproducible, queryable and
  maintainable through competency questions, integrity tests, deployment
  scripts and backup--restore validation.
\end{description}

The objectives consolidate the eleven operational goals used during the
sprints. Table~\ref{tab:goal-list} keeps the complete goal list visible.

\begin{table}[H]
\centering
\caption{Complete project-goal list}
\label{tab:goal-list}
\begin{tabular}{C{0.09\textwidth}L{0.48\textwidth}L{0.33\textwidth}}
\toprule
Goal & Description & Primary evidence \\
\midrule
G1 & Environmental Data Collection & Source register and collected files \\
G2 & Source Verification \& Selection & Exclusion register \\
G3 & Data Extraction & Extraction code and 0NF blocks \\
G4 & Data Cleaning \& Preprocessing & Data-quality log \\
G5 & ERD Development \& Integration & Final ERD and schema definition \\
G6 & Data Normalization to BCNF & Stage-wise outputs \\
G7 & Relational Schema Implementation & SQLite database \\
G8 & Data Loading \& ETL & Build scripts and normalized CSV files \\
G9 & Database Validation \& Integrity Checking & Automated tests \\
G10 & Query \& Demonstration Tools & Saved SQL and competency benchmark \\
G11 & Database Deployment \& Maintenance & Setup, backup and restore scripts \\
\bottomrule
\end{tabular}
\end{table}

Table~\ref{tab:objective-mapping} maps each problem to the objectives and report
evidence that address it.

\begin{table}[H]
\centering
\caption{Problem--objective--report mapping}
\label{tab:objective-mapping}
\begin{tabular}{C{0.10\textwidth}C{0.15\textwidth}L{0.61\textwidth}}
\toprule
Problem & Objective & Principal report evidence \\
\midrule
P1 & OB1, OB2 & Sections~\ref{sec:data}, \ref{sec:etl} and \ref{sec:norm} \\
P2 & OB1, OB2 & Sections~\ref{sec:data}, \ref{sec:etl} and \ref{sec:dm} \\
P3 & OB1, OB3 & Sections~\ref{sec:rga}, \ref{sec:deploy} and \ref{sec:findings} \\
\bottomrule
\end{tabular}
\end{table}

Figure~\ref{fig:problem-resolution} summarizes how the reported problems move
from a source condition to a recorded response and a verifiable outcome.

\begin{figure}[H]
\centering
\begin{tikzpicture}[
  problem/.style={draw=cublue,fill=cublue!8,rounded corners=2pt,
    text width=3.4cm,minimum height=1.25cm,align=center,font=\scriptsize},
  response/.style={draw=cuteal,fill=cuteal!8,rounded corners=2pt,
    text width=4.2cm,minimum height=1.25cm,align=center,font=\scriptsize},
  outcome/.style={draw=cuamber,fill=cuamber!8,rounded corners=2pt,
    text width=3.8cm,minimum height=1.25cm,align=center,font=\scriptsize},
  arr/.style={-{Stealth[length=2mm]},thick,draw=culine},
  node distance=0.45cm and 0.55cm
]
\node[problem] (p1) {P1: repeating columns, merged headers and no dependable source key};
\node[response,right=of p1] (r1) {Block extraction, unpivoting and staged normalization};
\node[outcome,right=of r1] (o1) {Atomic records and relations checked through BCNF};

\node[problem,below=of p1] (p2) {P2: inconsistent names and missing-value notation};
\node[response,right=of p2] (r2) {Canonical name maps, explicit missing values and quarantine rules};
\node[outcome,right=of r2] (o2) {Stable keys, reproducible joins and recorded exceptions};

\node[problem,below=of p2] (p3) {P3: mixed scope and grain, with unreliable join keys};
\node[response,right=of p3] (r3) {Documented source selection, grain-specific decomposition and validation};
\node[outcome,right=of r3] (o3) {Approved schema, integrity checks and repeatable SQLite build};

\draw[arr] (p1)--(r1); \draw[arr] (r1)--(o1);
\draw[arr] (p2)--(r2); \draw[arr] (r2)--(o2);
\draw[arr] (p3)--(r3); \draw[arr] (r3)--(o3);
\end{tikzpicture}
\caption{Problems identified, project responses and verified outcomes.}
\label{fig:problem-resolution}
\end{figure}

\subsection{Sector Description and Data Products}
\label{subsec:sdp}

The environment sector contains measurement families whose subjects and time
grains cannot be placed in one undifferentiated table. The implemented data
products are separated as follows.

\begin{description}[leftmargin=0.6cm,style=nextline]
  \item[Monthly climate records] Maximum and minimum temperature, humidity,
  rainfall, maximum and minimum wind speed with direction and thunderstorm and
  lightning observations are organized by station, year and month.
  \item[Daily climate records] Sunshine hours use station and calendar-day
  grain. Solar radiation adds a sample number because several readings can
  occur on the same day.
  \item[Hydrology and water quality] Rivers and monitoring stations support pH,
  dissolved oxygen, biochemical oxygen demand, chemical oxygen demand,
  salinity and marine-debris observations by station and year.
  \item[Forest area] Protected, reserved, unclassed, acquired and total forest
  areas are organized by district and fiscal period.
  \item[Industry] Establishment-size and industry-category observations retain
  their quantity, percentage and fiscal reporting period.
  \item[Reference and time data] Stations, districts, rivers, monitoring
  points, categories and calendar or fiscal periods provide consistent keys for
  the measurement relations.
\end{description}

\subsection{Scope, Sources and Definitions}
\label{subsec:sd}

Source discovery covers Bangladeshi and international organisations. The final
database uses nine approved structured files from Bangladesh Bureau of
Statistics (BBS), Bangladesh Meteorological Department (BMD), Bangladesh Rice
Research Institute (BRRI) and Bangladesh Water Development Board (BWDB)
\cite{bbs_env_ts,bmd_temperature,brri_daily,bwdb_rivers}. The wider collection
is retained as evidence even when a file cannot support the approved model.

Table~\ref{tab:organisations} adapts the implemented-source profile to the
verified Drive inventory. Collected files are source or data publications; the
loaded count excludes support files and derived page images.

\begin{table}[H]
\centering
\caption{Implemented source profile within the collected file inventory}
\label{tab:organisations}
\begin{tabular}{L{0.31\textwidth}r r L{0.38\textwidth}}
\toprule
Organisation & Collected & Loaded & Implemented material \\
\midrule
Bangladesh Bureau of Statistics (BBS) & 2 & 1 & Climate, water quality,
forest and industrial statistics. \\
Bangladesh Meteorological Department (BMD) & 11 & 2 & Station-based
temperature and sunshine publications. \\
Bangladesh Rice Research Institute (BRRI) & 9 & 5 & Daily temperature,
rainfall, humidity and sunshine observations. \\
Bangladesh Water Development Board (BWDB) & 4 & 1 & River names and
river-reference information. \\
\bottomrule
\end{tabular}
\end{table}

Discovery also covers the Department of Environment, Ministry of Environment,
Forest and Climate Change, Humanitarian Data Exchange, UNESCO, World Health
Organization, World Bank, United Nations and other publishers. Their reviewed
files use incompatible subjects or identifiers, lack a dependable join key or
consist of narrative material rather than repeatable records. They remain in
the source review without being represented as loaded database facts
\cite{hdx_bangladesh}.

In this report, a \textbf{block} means one rectangular region of published
data that shares one header structure. A sheet may hold several blocks. BBS
sheet T1.14, for example, publishes radiation for several stations as separate
tables stacked vertically and each becomes a separate block.

\subsection{Organization of the Report}
\label{subsec:org}

Sections~\ref{sec:projectmanagement}--\ref{sec:rga} cover management, Scrum and
sprints; Sections~\ref{sec:data}--\ref{sec:dm} cover sources, ETL,
normalization and modelling. Deployment, collaboration, findings and pedagogy
are followed by the evidence-based resource comparison. The conclusion closes
the numbered discussion, with the bibliography as the final report element.
